\documentclass[11pt,a4paper]{article}
\usepackage[utf8]{inputenc}
\usepackage[T1]{fontenc}
\usepackage[portuguese]{babel}
\usepackage{xcolor}
\usepackage{hyperref}
\usepackage{geometry}
\usepackage{tcolorbox}

\geometry{margin=2.5cm}

\definecolor{primary}{RGB}{41, 128, 185}
\definecolor{secondary}{RGB}{44, 62, 80}
\definecolor{codebg}{RGB}{240, 240, 240}

\hypersetup{colorlinks=true, linkcolor=primary, urlcolor=primary}

\title{\color{secondary}\Huge\textbf{Exercício 4.4: Seu canário}}
\author{\color{primary}DevOps com Kubernetes -- 2026}
\date{}

\begin{document}
\maketitle

\section*{\color{primary}Instruções do Exercício}
\begin{tcolorbox}[colback=blue!5!white,colframe=blue!75!black,title=Exercício 4.4: Seu canário]
Crie um \texttt{AnalysisTemplate} (utilizado pelo Argo Rollouts) para a aplicação Ping-pong que acompanhará e medirá o uso de CPU de todos os contêineres do seu \textit{namespace}.

A lógica de análise deve verificar se a taxa de uso da CPU (a soma) subir acima de um valor predefinido num intervalo de 5 minutos. Se o valor for excedido, o \textit{rollout} deve ser considerado falho e a atualização deve ser revertida (\textit{rollback}) automaticamente.
\end{tcolorbox}

\end{document}
